\section{Discussion and Limitations}
\paragraph{A useful front end, not complete reference understanding.}
The primary contribution is an operating point: a single compact image model can return a spatial proposal on glasses hardware. The gap between hand and target accuracy identifies the remaining difficulty. Detecting a plausible hand and fingertip is much easier than deciding which object is intended and drawing its full extent. This makes a local proposal stage a plausible component of an assistant, but its value depends on downstream tolerance for incorrect crops. Reductions in recognition cost, communication or task completion time were not measured and remain application-level implications.

The distinction between a proposal and an action is especially important for large-gap targets. Even at higher resolution, approximately two thirds of that subset fail \pacc{0.5}. A system should not treat a confident-looking box as permission for an irreversible action. Visual or verbal confirmation is a reasonable design implication, but no interaction outcome is claimed by this model evaluation.

\paragraph{Structure helps selectively.}
The ablations argue for evaluating mechanisms where they are intended to help, while still exposing the common headline metric. Instance conditioning improves stricter localisation under a common recipe; distractor supervision improves the specified interference condition; consistency losses improve a direction target without a demonstrated box benefit. The ray-branch result also prevents the architecture from being justified solely by its intuitive diagram. Learning with geometric channels is not equivalent to learning directional reasoning. A smaller A3 variant is a credible deployment candidate, but the prototype measurements and three-seed evidence belong to the full model, so we do not silently substitute the ablation as the reported system.

\paragraph{Data and generalisation.}
All quantitative model conclusions come from one released dataset, with unknown per-image provenance and incomplete demographic coverage. Its predominance of fingertip-overlap images limits what aggregate accuracy says about general pointing. The image gap is only a projected relation, and the square resize changes the geometry of non-square camera frames. Real glasses footage may differ in aspect ratio, motion blur, illumination, pose and camera--eye offset. Synthetic distractors cannot resolve these domain shifts; T5 specifically addresses real multi-hand behaviour, without turning a small check into a population-level generalisation claim.

The dataset also lacks the two conditional-negative states needed to establish general pointing and target-presence detection. A visible non-pointing hand may still trigger the cascade. Our localisation metrics do not hide these images through predicted gating, but the missing states limit deployment claims. The model predicts an image-plane target box and target-centre direction, not depth, anatomical pose or temporally coherent intent.

\paragraph{Evidence strength and reproducibility.}
Three seeds demonstrate stability of the principal full-model score, not invariance of every subgroup. Ablations and resolution variants have one seed, different stopping points and, in some cases, bundled interventions. Paired image intervals do not remove these limitations, especially for the 209 large-gap and 309 small-target examples. Published comparisons are useful context but cannot support claims of controlled superiority. The repository provides preprocessing, model/training code, variant configurations, checkpoints and per-image predictions; the accompanying audit script reconstructs the selected numerical summaries from these outputs.

\paragraph{Privacy and data practice.}
Local inference can keep the spatial-selection step on the glasses, but does not by itself make a head-worn camera private. A target crop may still contain bystanders, screens or identifying content. Consent, visible recording indication, retention limits and downstream data handling remain application responsibilities. EgoPoint is used under the authors' research-access terms; source images require separate access. The preprocessing code and evaluation outputs support reconstruction without making unrestricted redistribution of the source dataset part of the contribution. Consent and handling of new real frames are included in the T5 block.

\section{Conclusion}
SpatialNet demonstrates image-only pointing-target localisation with 722K parameters and reported execution on a smart-glasses prototype. Its value is a compact spatial front end whose strengths and boundaries are measured: instance conditioning improves precise boxes, distractor supervision improves synthetic interference robustness, and image detail yields the largest tested accuracy gain. Explicit ray conditioning does not establish a target-localisation benefit, and high hand accuracy does not resolve separated or ambiguous targets. The resulting evidence supports a practical design direction for local visual selection, while keeping real multi-hand validity and full interaction behaviour distinct from benchmark success.
